\documentclass[11pt]{scrartcl}

\usepackage[stixtwo]{fontsetup}
\usepackage{amsmath}

\begin{document}
\begin{enumerate}
	\item \[
		      a_n
		      = \frac{\left(\sqrt{n^2 + 2n} - n\right)\left(\sqrt{n^2 + 2n} + n\right)}{\sqrt{n^2 + 2n} + n}
		      = \frac{2n}{\sqrt{n^2 + 2n} + n}
		      = \frac{2}{\sqrt{1 + \frac{2}{n}} + 1}
		      \to 1
	      \]

	\item Minden $n \ge 1$ esetén
	      \(
	      \frac{n^2}{6} = \frac{n^3}{6n} \le \frac{n^3 + 1}{n + 5} \le \frac{2n^3}{n} = 2n^2
	      \),
	      így
	      \[
		      1
		      \leftarrow \sqrt[n]{\frac{1}{6}} \left(\sqrt[n]{n}\right)^2
		      \le b_n
		      \le \sqrt[n]{2} \left(\sqrt[n]{n}\right)^2
		      \to 1,
	      \]
	      tehát a rendőrelv miatt $b_n \to 1$.

	\item \[
		      c_n
		      = \frac{n!}{5^n} \cdot
		      \frac{1 + \frac{3^n}{n!}}{1 + \frac{n^3}{5^n}}
		      \to \infty,
	      \]
	      mert $\frac{n!}{5^n} \to \infty$, a második tényező pedig $1$-hez
	      tart, hiszen $\frac{3^n}{n!} \to 0$ és $\frac{n^3}{5^n} \to 0$.
\end{enumerate}
\end{document}
